\ifdefined\gkver\else\def\gkver{student}\fi
\documentclass[\gkver]{gaokaozhenti}
\begin{document}

\chapter{2009年重庆理综}
%% paperType: 理综物理部分

\begin{enumerate}
\item
%% number: 14
%% paperName: 2009年重庆理综第 14 题 6 分
%% typeId: 1
%% type: 单选题
%% chapter: 气体性质
%% point: 热力学第一定律
%% method: 无
%% score: 6
%% degree: 650
%% duplicateId: 0
%% body:
密闭有空气的薄塑料瓶因降温而变扁，此过程中瓶内空气（不计分子势能）\xzanswer{D}
\fourchoices[answer=D]
{内能增大，放出热量}
{内能减小，吸收热量}
{内能增大，对外界做功}
{内能减小，外界对其做功}

%% answer: D
%% memo:
\memoanswer{
不计分子势能，空气内能由温度决定、随温度降低而减小，AC 均错；薄塑料瓶因降温而变扁、空气体积减小、外界压缩空气做功，D 对；空气内能减少、外界对空气做功，根据热力学第一定律可知空气向外界放热、B 错。
}

\item
%% number: 15
%% paperName: 2009年重庆理综第 15 题 6 分
%% typeId: 1
%% type: 单选题
%% chapter: 机械波
%% point: 波的图象
%% method: 无
%% score: 6
%% degree: 650
%% duplicateId: 0
%% body:
同一音叉发出的声波同时在水和空气中传播，某时刻的波形曲线如图所示，以下说法正确的是\xzanswer{A}
\onepicture[w=7cm]{figs/15.png}
\fourchoices[answer=A]
{声波在水中波长较大，b 是水中声波的波形曲线}
{声波在空气中波长较大，b 是空气中声波的波形曲线}
{水中质点振动频率较高，a 是水中声波的波形曲线}
{空气中质点振动频率较高，a 是空气中声波的波形曲线}

%% answer: A
%% memo:
\memoanswer{
同一音叉发出的声波，声源相同，频率 $f$ 相同、周期 $T$ 相同（CD 均错）；又声波在水中传播的速度比在空气中快、速度 $v$ 大，根据波长 $\lambda=vT$ 可知声波在水中波长较大；由题图波形曲线可知 b 比 a 的波长长，b 是水中声波的波形曲线，A 对、B 错。
}

\item
%% number: 16
%% paperName: 2009年重庆理综第 16 题 6 分
%% typeId: 1
%% type: 单选题
%% chapter: 原子和原子核
%% point: 核聚变
%% method: 无
%% score: 6
%% degree: 650
%% duplicateId: 0
%% body:
某科学家提出年轻热星体中核聚变的一种理论，其中的两个核反应方程为\xzanswer{B}

$\ce{^1_1H + ^12_6C -> ^13_7N + Q_1}$ \qquad $\ce{^1_1H + ^15_7N -> ^12_6C + X + Q_2}$

方程式中 $Q_{1}$、$Q_{2}$ 表示释放的能量，相关的原子核质量见下表：

\begin{center}
\begin{tabular}{|c|c|c|c|c|c|c|}
\hline
原子核 & $\ce{^1_1H}$ & $\ce{^3_2He}$ & $\ce{^4_2He}$ & $\ce{^12_6C}$ & $\ce{^13_7N}$ & $\ce{^15_7N}$ \\
\hline
质量/$\Uu$ & 1.0078 & 3.0160 & 4.0026 & 12.0000 & 13.0057 & 15.0001 \\
\hline
\end{tabular}
\end{center}

\fourchoices[answer=B]
{X 是 $\ce{^3_2He}$，$Q_{2}>Q_{1}$}
{X 是 $\ce{^4_2He}$，$Q_{2}>Q_{1}$}
{X 是 $\ce{^3_2He}$，$Q_{2}<Q_{1}$}
{X 是 $\ce{^4_2He}$，$Q_{2}<Q_{1}$}

%% answer: B
%% memo:
\memoanswer{
$\ce{^1_1H + ^12_6C -> ^13_7N}$ 中质量亏损为 $\Delta m_{1}=1.0078\Uu+12.0000\Uu-13.0057\Uu=0.0021\Uu$。

根据电荷数守恒和质量数守恒可知 $\ce{^1_1H + ^15_7N -> ^12_6C + X}$ 中 X 的电荷数为 2、质量数为 4，质量亏损为 $\Delta m_{2}=1.0078\Uu+15.0001\Uu-12.0000\Uu-4.0026\Uu=0.0053\Uu$，根据爱因斯坦的质能方程可知 $Q_{1}=\Delta m_{1}c^{2}$、$Q_{2}=\Delta m_{2}c^{2}$，则 $Q_{1}<Q_{2}$。
}

\item
%% number: 17
%% paperName: 2009年重庆理综第 17 题 6 分
%% typeId: 1
%% type: 单选题
%% chapter: 曲线运动
%% point: 线速度角速度
%% method: 无
%% score: 6
%% degree: 650
%% duplicateId: 0
%% body:
据报道，“嫦娥一号”和“嫦娥二号”绕月飞行器的圆形轨道距月球表面分别约为 $200\Ukm$ 和 $100\Ukm$，运动速率分别为 $v_{1}$ 和 $v_{2}$。那么 $v_{1}$ 和 $v_{2}$ 的比值为（月球半径取 $1700\Ukm$）\xzanswer{C}
\fourchoices[answer=C]
{$\frac{19}{18}$}
{$\sqrt{\frac{19}{18}}$}
{$\sqrt{\frac{18}{19}}$}
{$\frac{18}{19}$}

%% answer: C
%% memo:
\memoanswer{
“嫦娥一号”和“嫦娥二号”绕月作圆周运动，由万有引力提供向心力有 $\frac{GMm}{R^{2}}=\frac{mv^{2}}{R}$，可得 $v=\sqrt{\frac{GM}{R}}$（$M$ 为月球质量），它们的轨道半径分别为 $R_{1}=1900\Ukm$、$R_{2}=1800\Ukm$，则 $v_{1}:v_{2}=\sqrt{\frac{R_{2}}{R_{1}}}=\sqrt{\frac{18}{19}}$。
}

\item
%% number: 18
%% paperName: 2009年重庆理综第 18 题 6 分
%% typeId: 1
%% type: 单选题
%% chapter: 电路
%% point: 动态分析
%% method: 无
%% score: 6
%% degree: 450
%% duplicateId: 0
%% body:
某实物投影机有 10 个相同的强光灯 $L_{1}\sim L_{10}$（$24\UV/200\UW$）和 10 个相同的指示灯 $X_{1}\sim X_{10}$（$220\UV/2\UW$），将其连接在 $220\UV$ 交流电源上，电路如图所示，若工作一段时间后，$L_{2}$ 灯丝烧断，则\xzanswer{C}
\onepicture[w=7cm]{figs/18.png}
\fourchoices[answer=C]
{$X_{1}$ 的功率减小，$L_{1}$ 的功率增大}
{$X_{1}$ 的功率增大，$L_{1}$ 的功率增大}
{$X_{2}$ 功率增大，其它指示灯的功率减小}
{$X_{2}$ 功率减小，其它指示灯的功率增大}

%% answer: C
%% memo:
\memoanswer{
显然 $L_{1}$ 和 $X_{1}$ 并联、$L_{2}$ 和 $X_{2}$ 并联……然后它们再串联接在 $220\UV$ 交流电源上，$L_{2}$ 灯丝烧断，则总电阻变大、电路中电流 $I$ 减小，又 $L_{1}$ 和 $X_{1}$ 并联的电流分配关系不变，则 $X_{1}$ 和 $L_{1}$ 的电流都减小、功率都减小，同理可知除 $X_{2}$ 和 $L_{2}$ 外各灯功率都减小，A、B 均错；由于 $I$ 减小，各并联部分的电压都减小，交流电源电压不变，则 $X_{2}$ 上电压增大，根据 $P=\frac{U^{2}}{R}$ 可知 $X_{2}$ 的功率变大，C 对、D 错。
}

\item
%% number: 19
%% paperName: 2009年重庆理综第 19 题 6 分
%% typeId: 1
%% type: 单选题
%% chapter: 磁场
%% point: 左手定则
%% method: 无
%% score: 6
%% degree: 650
%% duplicateId: 0
%% body:
在如图所示电路中，电池均相同，当电键 S 分别置于 a、b 两处时，导线 $MM'$ 与 $NN'$ 之间的安培力的大小为 $f_{a}$、$f_{b}$，判断这两段导线\xzanswer{D}
\onepicture[w=6.5cm]{figs/19.png}
\fourchoices[answer=D]
{相互吸引，$f_{a}>f_{b}$}
{相互排斥，$f_{a}>f_{b}$}
{相互吸引，$f_{a}<f_{b}$}
{相互排斥，$f_{a}<f_{b}$}

%% answer: D
%% memo:
\memoanswer{
电键 S 分别置于 a、b 两处时，电源分别为一节干电池、两节干电池，而电路中灯泡电阻不变，则电路中电流 $I_{a}<I_{b}$，$MM'$ 在 $NN'$ 处的磁感应强度 $B_{a}<B_{b}$，应用安培力公式 $F=BIL$ 可知 $f_{a}<f_{b}$，又 $MM'$ 与 $NN'$ 中电流方向相反，则相互排斥。
}

\item
%% number: 20
%% paperName: 2009年重庆理综第 20 题 6 分
%% typeId: 1
%% type: 单选题
%% chapter: 电磁感应
%% point: 切割磁感线电动势
%% method: 无
%% score: 6
%% degree: 450
%% duplicateId: 0
%% body:
如图所示为一种早期发电机原理示意图，该发电机由固定的圆形线圈和一对用铁芯连接的圆柱形磁铁构成，两磁极相对于线圈平面对称，在磁极绕转轴匀速转动过程中，磁极中心在线圈平面上的投影沿圆弧 XOY 运动（O 是线圈中心），则\xzanswer{D}
\onepicture[w=5cm]{figs/20.png}
\fourchoices[answer=D]
{从 X 到 O，电流由 E 经 G 流向 F，先增大再减小}
{从 X 到 O，电流由 F 经 G 流向 E，先减小再增大}
{从 O 到 Y，电流由 F 经 G 流向 E，先减小再增大}
{从 O 到 Y，电流由 E 经 G 流向 F，先增大再减小}

%% answer: D
%% memo:
\memoanswer{
在磁极绕转轴从 X 到 O 匀速转动过程中，穿过线圈平面的磁通量向上增大，根据楞次定律可知线圈中产生顺时针方向的感应电流，电流由 F 经 G 流向 E，又导线切割磁感线产生感应电动势 $E_{\text{感}}=BLv$，导线处的磁感应强度先增后减可知感应电动势先增加后减小、则电流先增大再减小，AB 均错；

在磁极绕转轴从 O 到 Y 匀速转动的过程中，穿过线圈平面的磁通量向上减小，根据楞次定律可知线圈中产生逆时针方向的感应电流，电流由 E 经 G 流向 F，又导线切割磁感线产生感应电动势 $E_{\text{感}}=BLv$，导线处的磁感应强度先增后减可知感应电动势先增加后减小、则电流先增大再减小，C 错、D 对。
}

\item
%% number: 21
%% paperName: 2009年重庆理综第 21 题 6 分
%% typeId: 1
%% type: 单选题
%% chapter: 光学
%% point: 光的干涉
%% method: 无
%% score: 6
%% degree: 450
%% duplicateId: 0
%% body:
用 a、b、c、d 表示四种单色光，若

①a、b 从同种玻璃射向空气，a 的临界角小于 b 的临界角；

②用 b、c 和 d 在相同条件下分别做双缝干涉实验，c 的条纹间距最大；

③用 b、d 照射某金属表面，只有 b 能使其发射电子。

则可推断 a、b、c、d 可能分别是\xzanswer{A}
\fourchoices[answer=A]
{紫光、蓝光、红光、橙光}
{蓝光、紫光、红光、橙光}
{紫光、蓝光、橙光、红光}
{紫光、橙光、红光、蓝光}

%% answer: A
%% memo:
\memoanswer{
根据临界角 $C$、折射率 $n=\frac{1}{\sin C}$，由条件①可知 $n_{a}>n_{b}$，根据色散规律可知 a 的频率大于 b；根据双缝干涉条纹间距 $\Delta X=\frac{L}{d}\lambda$，由条件②可知 b、c 和 d 中 c 的波长最长，再根据色散规律可知 bcd 中 c 的频率最小；每种金属都有对应的最小入射光频率，入射光频率越大、光电效应越容易发生，由条件③可知 b 和 d 中 b 的频率大，综合上述可知 a、b、c、d 的频率从大到小依次为 abdc，只有 A 选项中满足。
}

\item
%% number: 22
%% paperName: 2009年重庆理综第 22 题 12 分
%% typeId: 4
%% type: 实验题
%% chapter: 电路
%% point: 电路实验
%% method: 无
%% score: 12
%% degree: 450
%% duplicateId: 0
%% body:
\begin{enumerate}
	\item
	某同学在探究影响单摆周期的因素时有如下操作，请判断是否恰当（填“是”或“否”）。

	①把单摆从平衡位置拉开约 $5^{\circ}$ 释放；\tkanswer[2.2cm]{}

	②在摆球经过最低点时启动秒表计时；\tkanswer[2.2cm]{}

	③把秒表记录摆球一次全振动的时间作为周期。\tkanswer[2.2cm]{}

	该同学改进测量方法后，得到的部分测量数据见表。

	\begin{center}
	\begin{tabular}{|c|c|c|c|}
	\hline
	数据组编号 & 摆长/$\Umm$ & 摆球质量/$\Ug$ & 周期/$\Us$ \\
	\hline
	1 & 999.3 & 32.2 & 2.0 \\
	\hline
	2 & 999.3 & 16.5 & 2.0 \\
	\hline
	3 & 799.2 & 32.2 & 1.8 \\
	\hline
	4 & 799.2 & 16.5 & 1.8 \\
	\hline
	5 & 501.1 & 32.2 & 1.4 \\
	\hline
	\end{tabular}
	\end{center}

	用螺旋测微器测量其中一个摆球直径的示数如图所示。该球的直径为\tkanswer[2.2cm]{} $\Umm$。根据表中数据可以初步判断单摆周期随\tkanswer[2.2cm]{}的增大而增大。

	\onepicture[w=3.5cm, align=right]{figs/22.png}

	\item
	硅光电池是一种可将光能转换为电能的器件。某同学用图 \ref{2009重庆理综22a} 所示电路探究硅光电池的路端电压 $U$ 与总电流 $I$ 的关系。图中 $R_{0}$ 为已知定值电阻，电压表视为理想电压表。

	\twopicture[labelA=2009重庆理综22a, labelB=2009重庆理综22b, numA=1, numB=2, numstyle=arabic, labelprefix={图}, wA=4.6cm, wB=8cm, align=right]
	{figs/22a.png}
	{figs/22b.png}

	①请根据图 \ref{2009重庆理综22a}，用笔画线代替导线将图 \ref{2009重庆理综22b} 中的实验器材连接成实验电路。

	②若电压表 $V_{2}$ 的读数为 $U_{0}$，则 $I=$\tkanswer[2.2cm]{} $\UmA$；

	③实验一：用一定强度的光照射硅光电池，调节滑动变阻器，通过测量得到该电池的 $U-I$ 曲线 a。见图 \ref{2009重庆理综22c}，由此可知电池内阻\tkanswer[2.2cm]{}（填“是”或“不是”）常数，短路电流为\tkanswer[2.2cm]{} $\UmA$，电动势为\tkanswer[2.2cm]{} $\UV$。

	\onepicture[num=3, numstyle=arabic, labelprefix={图}, label=2009重庆理综22c, w=6cm, align=right]{figs/22c.png}

	④实验二：减小实验一中光的强度，重复实验，测得 $U-I$ 曲线 b，见图 \ref{2009重庆理综22c}。当滑动变阻器的电阻为某值时，若实验一中的路端电压为 $1.5\UV$。则实验二中外电路消耗的电功率为\tkanswer[2.2cm]{} $\UmW$（计算结果保留两位有效数字）。
\end{enumerate}

%% answer:
\jdanswer{
\begin{enumerate}
	\item
	是；是；否；$20.685$（$20.683\sim20.687$），摆长。
	\item
	①如图；②$\frac{U_{0}}{R_{0}}$；③不是，$0.295$（$0.293\sim0.297$），$2.67$（$2.64\sim2.70$）；④$0.065$（$0.060\sim0.070$）。
\end{enumerate}
}
%% memo:
\memoanswer{
（1）单摆作简谐运动要求摆角小，单摆从平衡位置拉开约 $5^{\circ}$ 释放满足此条件；因为最低点位置固定、容易观察，所以在最低点启动秒表计时；摆球一次全振动的时间太短、不易读准、误差大，应测多个周期的时间求平均值；表中数据可以初步判断单摆周期随摆长的增大而增大。

（2）见右图；②根据欧姆定律可知 $I=\frac{U_{0}}{R_{0}}$；③路端电压 $U=E-Ir$，若 $r$ 为常数、则 $U-I$ 图为一条不过原点的直线，由曲线 a 可知电池内阻不是常数；当 $U=0$ 时的电流为短路电流、约为 $295\UuA=0.295\UmA$；当电流 $I=0$ 时路端电压等于电源电动势 $E$、约为 $2.67\UV$；④实验一中的路端电压为 $U_{1}=1.5\UV$ 时电路中电流为 $I_{1}=0.21\UmA$，连接 a 中点（$0.21\UmA$、$1.5\UV$）和坐标原点，此直线为此时对应滑动变阻器阻值的外电路电阻（定值电阻）的 $U-I$ 图，和图线 b 的交点为实验二中的路端电压和电路电流，如右图，电流和电压分别为 $I=97\UuA$、$U=0.7\UV$，则外电路消耗功率为 $P=UI=0.068\UmW$。

\onepicture[showlabel=false, w=6cm, align=right]{figs/22_an1.png}

\onepicture[showlabel=false, w=6.5cm, align=right]{figs/22_an2.png}
}

\item
%% number: 23
%% paperName: 2009年重庆理综第 23 题 16 分
%% typeId: 6
%% type: 计算题
%% chapter: 功和能
%% point: 动能定理
%% method: 无
%% score: 16
%% degree: 450
%% duplicateId: 0
%% body:
2009 年中国女子冰壶队首次获得了世界锦标赛冠军，这引起了人们对冰壶运动的关注。冰壶在水平冰面上的一次滑行可简化为如下过程：如图所示，运动员将静止于 O 点的冰壶（视为质点）沿直线 $OO'$ 推到 A 点放手，此后冰壶沿 $AO'$ 滑行，最后停于 C 点。已知冰面与各冰壶间的动摩擦因数为 $\mu$，冰壶质量为 $m$，$AC=L$，$CO'=r$，重力加速度为 $g$。
\begin{enumerate}
	\item
	求冰壶在 A 点的速率；
	\item
	求冰壶从 O 点到 A 点的运动过程中受到的冲量大小；
	\item
	若将 $BO'$ 段冰面与冰壶间的动摩擦因数减小为 $0.8\mu$，原只能滑到 C 点的冰壶能停于 $O'$ 点，求 A 点与 B 点之间的距离。
\end{enumerate}
\onepicture[w=9cm, align=right]{figs/23.png}

%% answer:
\jdanswer{
\begin{enumerate}
	\item
	$\sqrt{2\mu gL}$
	\item
	$I=m\sqrt{2\mu gL}$
	\item
	$L-4r$
\end{enumerate}
}
%% memo:
\memoanswer{
（1）对冰壶，从 A 点放手到停止于 C 点，设在 A 点时的速度为 $v_{1}$，应用动能定理有 $-\mu mgL=0-\frac{1}{2}mv_{1}^{2}$，解得 $v_{1}=\sqrt{2\mu gL}$；

（2）对冰壶，从 O 到 A，设冰壶受到的冲量为 $I$，应用动量定理有 $I=mv_{1}-0$，解得 $I=m\sqrt{2\mu gL}$；

（3）设 AB 之间距离为 $s$，对冰壶，从 A 到 $O'$ 的过程，应用动能定理有 $-\mu mgs-0.8\mu mg(L+r-s)=0-\frac{1}{2}mv_{1}^{2}$，解得 $s=L-4r$。
}

\item
%% number: 24
%% paperName: 2009年重庆理综第 24 题 18 分
%% typeId: 6
%% type: 计算题
%% chapter: 运动和力
%% point: 动量和能量
%% method: 无
%% score: 18
%% degree: 450
%% duplicateId: 0
%% body:
探究某种笔的弹跳问题时，把笔分为轻质弹簧、内芯和外壳三部分，其中内芯和外壳质量分别为 $m$ 和 $4m$。笔的弹跳过程分为三个阶段：

①把笔竖直倒立于水平硬桌面，下压外壳使其下端接触桌面（见图 a）；

②由静止释放，外壳竖直上升至下端距桌面高度为 $h_{1}$ 时，与静止的内芯碰撞（见图 b）；

③碰后，内芯与外壳以共同的速度一起上升到外壳下端距桌面最大高度为 $h_{2}$ 处（见图 c）。

设内芯与外壳的撞击力远大于笔所受重力、不计摩擦与空气阻力，重力加速度为 $g$。求：
\begin{enumerate}
	\item
	外壳与内芯碰撞后瞬间的共同速度大小；
	\item
	从外壳离开桌面到碰撞前瞬间，弹簧做的功；
	\item
	从外壳下端离开桌面到上升至 $h_{2}$ 处，笔损失的机械能。
\end{enumerate}
\onepicture[w=5cm, align=right]{figs/24.png}

%% answer:
\jdanswer{
\begin{enumerate}
	\item
	$\sqrt{2g(h_{2}-h_{1})}$
	\item
	$\frac{25h_{2}-9h_{1}}{4}mg$
	\item
	$E_{\text{损}}=\frac{5}{4}mg(h_{2}-h_{1})$
\end{enumerate}
}
%% memo:
\memoanswer{
设外壳上升高度 $h_{1}$ 时速度为 $V_{1}$，外壳与内芯碰撞后瞬间的共同速度大小为 $V_{2}$，

（1）对外壳和内芯，从撞后达到共同速度到上升至 $h_{2}$ 处，应用动能定理有 $(4m+m)g(h_{2}-h_{1})=\frac{1}{2}(4m+m)V_{2}^{2}$，解得 $V_{2}=\sqrt{2g(h_{2}-h_{1})}$；

（2）外壳和内芯，碰撞过程瞬间动量守恒，有 $4mV_{1}=(4m+m)V_{2}$，解得 $V_{1}=\frac{5}{4}\sqrt{2g(h_{2}-h_{1})}$，设从外壳离开桌面到碰撞前瞬间弹簧做功为 $W$，在此过程中，对外壳应用动能定理有 $W-4mgh_{1}=\frac{1}{2}(4m)V_{1}^{2}$，解得 $W=\frac{25h_{2}-9h_{1}}{4}mg$；

（3）由于外壳和内芯达到共同速度后上升高度 $h_{2}$ 的过程，机械能守恒，只是在外壳和内芯碰撞过程有能量损失，损失的能量为 $E_{\text{损}}=\frac{1}{2}(4m)V_{1}^{2}-\frac{1}{2}(4m+m)V_{2}^{2}$，联立解得 $E_{\text{损}}=\frac{5}{4}mg(h_{2}-h_{1})$。
}

\item
%% number: 25
%% paperName: 2009年重庆理综第 25 题 19 分
%% typeId: 6
%% type: 计算题
%% chapter: 磁场
%% point: 带电粒子在复合场中的运动
%% method: 分情况讨论
%% score: 19
%% degree: 250
%% duplicateId: 0
%% body:
如图所示，离子源 A 产生的初速为零、带电量均为 $e$、质量不同的正离子被电压为 $U_{0}$ 的加速电场加速后匀速通过准直管，垂直射入匀强偏转电场，偏转后通过极板 HM 上的小孔 S 离开电场，经过一段匀速直线运动，垂直于边界 MN 进入磁感应强度为 $B$ 的匀强磁场。已知 $HO=d$，$HS=2d$，$\angle MNQ=90^{\circ}$。（忽略粒子所受重力）
\begin{enumerate}
	\item
	求偏转电场场强 $E_{0}$ 的大小以及 HM 与 MN 的夹角 $\varphi$；
	\item
	求质量为 $m$ 的离子在磁场中做圆周运动的半径；
	\item
	若质量为 $4m$ 的离子垂直打在 NQ 的中点 $S_{1}$ 处，质量为 $16m$ 的离子打在 $S_{2}$ 处。求 $S_{1}$ 和 $S_{2}$ 之间的距离以及能打在 NQ 上的正离子的质量范围。
\end{enumerate}
\onepicture[w=5.5cm, align=right]{figs/25.png}

%% answer:
\jdanswer{
\begin{enumerate}
	\item
	$E_{0}=\frac{U_{0}}{d}$，$\varphi=45^{\circ}$
	\item
	$R=2\sqrt{\frac{mU_{0}}{eB^{2}}}$
	\item
	$\Delta S=4(\sqrt{3}-1)\sqrt{\frac{mU_{0}}{eB^{2}}}$，$m<m_{x}<25m$
\end{enumerate}
}
%% memo:
\memoanswer{
（1）正离子被电压为 $U_{0}$ 的加速电场加速后速度设为 $v_{1}$，对正离子，应用动能定理有 $eU_{0}=\frac{1}{2}mv_{1}^{2}$，正离子垂直射入匀强偏转电场，作类平抛运动，受到电场力 $F=qE_{0}$、产生的加速度为 $a=\frac{F}{m}$，即 $a=\frac{qE_{0}}{m}$，垂直电场方向匀速运动，有 $2d=v_{1}t$，沿场强方向 $Y=\frac{1}{2}at^{2}$，联立解得 $E_{0}=\frac{U_{0}}{d}$，又 $\tan\varphi=\frac{v_{1}}{at}$，解得 $\varphi=45^{\circ}$；

（2）正离子进入磁场时的速度大小为 $v_{2}=\sqrt{v_{1}^{2}+(at)^{2}}$，正离子在匀强磁场中作匀速圆周运动，由洛伦兹力提供向心力 $qv_{2}B=\frac{mv_{2}^{2}}{R}$，解得离子在磁场中做圆周运动的半径 $R=2\sqrt{\frac{mU_{0}}{eB^{2}}}$；

（3）根据 $R=2\sqrt{\frac{mU_{0}}{eB^{2}}}$ 可知，质量为 $4m$ 的离子在磁场中的运动打在 $S_{1}$，运动半径为 $R_{1}=2\sqrt{\frac{4mU_{0}}{eB^{2}}}$，质量为 $16m$ 的离子在磁场中的运动打在 $S_{2}$，运动半径为 $R_{2}=2\sqrt{\frac{16mU_{0}}{eB^{2}}}$，又 $ON=R_{2}-R_{1}$，由几何关系可知 $S_{1}$ 和 $S_{2}$ 之间的距离 $\Delta S=\sqrt{R_{2}^{2}-ON^{2}}-R_{1}$，联立解得 $\Delta S=4(\sqrt{3}-1)\sqrt{\frac{mU_{0}}{eB^{2}}}$；由 $R'^{2}=(2R_{1})^{2}+(R'-R_{1})^{2}$ 解得 $R'=\frac{5}{2}R_{1}$，再根据 $\frac{1}{2}R_{1}<R<\frac{5}{2}R_{1}$，解得 $m<m_{x}<25m$。

\onepicture[showlabel=false, w=5cm, align=right]{figs/25_an1.png}

\onepicture[showlabel=false, w=3.5cm, align=right]{figs/25_an2.png}
}

\end{enumerate}

\end{document}
